\documentclass{article}
\usepackage[margin=28mm]{geometry}
\usepackage{amsmath,amssymb}
\usepackage{xurl}
\usepackage[hidelinks]{hyperref}
% Shared commands for notes in docs/paper-gaps/.

\DeclareUnicodeCharacter{2080}{\ifmmode{}_{0}\else\textsubscript{0}\fi}
\DeclareUnicodeCharacter{2081}{\ifmmode{}_{1}\else\textsubscript{1}\fi}
\DeclareUnicodeCharacter{2082}{\ifmmode{}_{2}\else\textsubscript{2}\fi}
\DeclareUnicodeCharacter{2099}{\ifmmode{}_{n}\else\textsubscript{n}\fi}

\newcommand{\Lean}{\textsc{Lean}}
\ExplSyntaxOn
\NewDocumentCommand{\leanid}{m}{
  \ifmmode
    \textsf{\detokenize{#1}}
  \else
    \group_begin:
    \urlstyle{sf}
    \tl_set:Nn \l_tmpa_tl {#1}
    \tl_replace_all:Nnn \l_tmpa_tl {\_} {_}
    \exp_args:NV \path \l_tmpa_tl
    \group_end:
  \fi
}
\ExplSyntaxOff
\providecommand{\tr}{\operatorname{tr}}
\newcommand{\tnleanIssueBase}{https://github.com/LionSR/TNLean/issues/}
\newcommand{\tnleanPrBase}{https://github.com/LionSR/TNLean/pull/}
\newcommand{\tnleanDocBase}{https://sirui-lu.com/TNLean/docs/}
\newcommand{\ghissue}[1]{\href{\tnleanIssueBase#1}{GitHub issue \##1}}
\newcommand{\ghpr}[1]{\href{\tnleanPrBase#1}{PR \##1}}
\newcommand{\doclink}[1]{\href{\tnleanDocBase#1.html}{\path{#1}}}

% Dirac notation and the identity operator, as in blueprint/src/macros/common.tex.
% \providecommand keeps note-local definitions authoritative.
\usepackage{dsfont}
\providecommand{\ket}[1]{|#1\rangle}
\providecommand{\bra}[1]{\langle #1|}
\providecommand{\braket}[2]{\langle #1 | #2 \rangle}
\providecommand{\ketbra}[2]{|#1\rangle\!\langle #2|}
\providecommand{\Id}{\mathds{1}}
\providecommand{\one}{\mathds{1}}
\providecommand{\C}{\mathbb{C}}
\providecommand{\MN}[1]{M_{#1}(\C)}
\providecommand{\MD}{M_D(\C)}

% Verdict marker: \gapnote{<kind>}{<status>}. Kind is the policy
% classification (clarification, local-correction, scope-restriction,
% unfaithful, false-source, open-gap); status is open, wip (elimination
% underway), resolved, or historical. Machine-read by the site index;
% prints a verdict line.
\newcommand{\gapnote}[2]{\par\noindent\textbf{Verdict:} #1 (#2).\par}

\providecommand{\gapnote}[2]{\par\noindent\textbf{Verdict:} #1 (#2).\par}
\providecommand{\leanid}[1]{\textsf{\detokenize{#1}}}
\title{The Boundary of a Cluster in Hastings' High-Temperature Expansion}
\author{}
\date{}
\begin{document}
\maketitle
\gapnote{local-correction}{resolved}

\section{The source statement}
Hastings, \emph{Solving gapped Hamiltonians locally}~\cite{Hastings2006Solving},
\href{https://arxiv.org/abs/cond-mat/0508554v3}{arXiv:cond-mat/0508554v3},
source \path{proj3new.tex}, lines 1130--1138, reads:
``Let $C$ be some set of sites $i$. Define $H_C=\sum_{i\in C}H_i$. Define
$B(C)$ to be equal to the set of all sites $j$ such that $j\neq i$ for any
$i\in C$ and such that $H_j$ acts on a site $i$ for some site $i\in C$. [\dots]
Then, if a given term in the power series expansion for $\exp(-\beta H)$
includes $H_i$ for every $i\in C$, but does not include $H_j$ for all $j\in
B(C)$, then this term in the power series expansion has a cluster which
includes exactly the sites acted on by $H_i$ for all $i\in C$.''
The terms are indexed by sites, and $H_i$ acts only on sites within the
interaction range $R$ of $i$ (lines 181--183). The operator $\rho(\beta,C)$
(lines 1174--1177) and the inclusion--exclusion sum
$\rho-\rho_1+\rho_2-\cdots$ of the remark at lines 1218--1224 rest on this
claim.

\section{The point at issue}
The claim fails for the boundary as defined. Take a chain with $R=1$, so that
$H_i$ acts on the sites $\{i-1,i,i+1\}$, and $C=\{1\}$. The terms acting on the
site $1$ are $H_0$, $H_1$ and $H_2$, so $B(C)=\{0,2\}$. The term $H_1H_3$ of
the power series contains $H_1$ and avoids $H_0$ and $H_2$, but $H_3$ acts on
$\{2,3,4\}$, which meets the sites $\{0,1,2\}$ of $H_1$ at the site $2$. The
active bonds of this term therefore join $\{0,\dots,4\}$ into one cluster, and
no cluster consists of exactly the sites of $H_1$. The same happens with
$H_i$ acting on $\{i,i+1\}$: then $B(\{1\})=\{0\}$, and $H_1H_2$ avoids it
while $H_2$ meets $H_1$ at the site $2$.

With this boundary, $\rho(\beta,C)$ counts terms in which $C$ is not a
cluster, and the inclusion--exclusion sum over families of large clusters
does not reproduce $\rho(\beta,l_{\mathrm{proj}})$.

\section{The correction}
The boundary is taken to be the set of terms $j\notin C$ whose support meets
the support of $H_C$:
\begin{align}
  \partial C=\Bigl\{j\notin C:\ \operatorname{supp}H_j\cap
  \bigcup_{i\in C}\operatorname{supp}H_i\neq\emptyset\Bigr\}.
  \label{eq:hastings05_cluster_boundary_overlap}
\end{align}
These are the neighbours of $C$ in the graph on the terms in which two
distinct terms are adjacent when their supports meet. For a set $C$ of terms
that is connected in this graph, a word containing every $H_i$ with $i\in C$
and no $H_j$ with $j\in\partial C$ has the terms of $C$ as one of its
clusters, and conversely. A set $C$ that is not connected splits into several
clusters, so the claim is used only for connected $C$, which are the
clusters of the expansion.

\section{The effect downstream}
The second inequality of line 1186, $n_C+n_{B(C)}\le Sn_C$, is unchanged in
form but $S$ changes meaning. With the boundary
(\ref{eq:hastings05_cluster_boundary_overlap}), $S$ must bound, for each term,
the number of terms whose support meets its support, together with the term
itself. When every support lies within $R$ of its centre, these are terms
whose centres lie within $2R$, so $S$ counts the sites within distance $2R$ of
a site rather than within $R$. The weight $y=(e^{J\beta}-1)e^{SJ\beta}$ of
equation (percb2), and with it the temperature $\beta_0$ below which
$xy\le1/2$ (line 1244), change accordingly. The percolation length
$\xi_{\mathrm{perc}}$ changes only through $y$ and remains of order the
interaction range $R$, as the source states (line 1250).

\section{Formalization}
The overlap graph and the boundary
(\ref{eq:hastings05_cluster_boundary_overlap}) are
\leanid{TNLean.Lattice.ClusterExpansion.overlapGraph} and
\leanid{TNLean.Lattice.ClusterExpansion.mem_boundary_overlapGraph}. Clusters of
a set of terms are
\leanid{TNLean.Lattice.ClusterExpansion.IsCluster}, and the
inclusion--exclusion identity is
\leanid{TNLean.Lattice.ClusterExpansion.hasSum_truncatedExpTerm}. The bounds on
$\rho(\beta,C)$ are stated for an arbitrary boundary, so they apply unchanged.
\bibliographystyle{unsrt}
\bibliography{references}
\end{document}
